\documentclass[11pt,a4paper]{article}

\usepackage[utf8]{inputenc}
\usepackage[T1]{fontenc}
\usepackage[french]{babel}
\usepackage{amsmath,amssymb,amsthm}
\usepackage{geometry}
\usepackage{booktabs}
\usepackage{enumitem}

\geometry{margin=2.5cm}

\newcommand{\R}{\mathbb{R}}
\newcommand{\X}{X}
\newcommand{\Y}{Y}
\newcommand{\ps}[2]{\left\langle #1,\, #2 \right\rangle}
\newcommand{\nr}[1]{\left\| #1 \right\|}
\newcommand{\dive}{\operatorname{div}}
\newcommand{\F}{\mathcal{F}}

\newtheorem*{remarque}{Remarque}

\title{TP2 -- Traitement d'images par approches variationnelles\\[2pt]
\large Réponses aux questions théoriques}
\author{Younes Bouhoreira}
\date{\today}

\begin{document}
\maketitle

\section*{Notations}

On travaille dans $\X = \R^{m\times n}$ muni du produit scalaire $\ps{u}{v}_\X = \sum_{x\in\Omega} u(x)v(x)$,
et dans $\Y = \X \times \X$ (champs de vecteurs) muni de $\ps{p}{q}_\Y = \sum_{x\in\Omega} p(x)\cdot q(x)$.
Le gradient discret $\nabla : \X \to \Y$ est linéaire, et la divergence discrète est définie par dualité :
\[
  \dive = -\nabla^{*}, \qquad \text{c'est-à-dire} \qquad
  \ps{\nabla u}{p}_\Y = -\ps{u}{\dive p}_\X \quad \forall u\in\X,\ \forall p \in \Y .
\]
C'est exactement la construction des fonctions \texttt{gradient} et \texttt{div} du notebook
(différences avant, conditions de Neumann). On note $\Delta = \dive \nabla$.

Pour montrer qu'une fonction $J$ est différentiable de gradient $\nabla J(u)$, on développe
$J(u+h) = J(u) + \ps{\nabla J(u)}{h}_\X + o(\nr{h})$.

\section*{Question 1 -- Convexité et gradients de $J_H$ et $J_{PM}$}

\subsection*{Fonction $J_H(u) = \frac12 \sum_{x} \nr{\nabla u(x)}^2 = \frac12 \nr{\nabla u}_\Y^2$}

\paragraph{Convexité.}
L'application $p \mapsto \frac12\nr{p}_\Y^2$ est convexe et $\nabla$ est linéaire.
Pour $u,v\in\X$ et $t\in[0,1]$, on a l'identité
\[
  \tfrac12\nr{t\,a + (1-t)\,b}^2 = t\,\tfrac12\nr{a}^2 + (1-t)\,\tfrac12\nr{b}^2 - \tfrac12\,t(1-t)\nr{a-b}^2 ,
\]
appliquée avec $a = \nabla u$, $b=\nabla v$ :
\[
  J_H(tu+(1-t)v) = t J_H(u) + (1-t) J_H(v) - \tfrac12\,t(1-t)\nr{\nabla(u-v)}_\Y^2 \le t J_H(u) + (1-t) J_H(v).
\]
Donc $J_H$ est \textbf{convexe}. Elle n'est \textbf{pas strictement convexe} : si $u - v = c$ est une
image constante, $\nabla(u-v)=0$ et on a égalité. Le noyau de $\nabla$ est l'ensemble $C$ des images constantes.

\paragraph{Gradient.}
\begin{align*}
  J_H(u+h) &= \tfrac12\nr{\nabla u + \nabla h}_\Y^2
  = J_H(u) + \ps{\nabla u}{\nabla h}_\Y + \tfrac12\nr{\nabla h}_\Y^2 \\
  &= J_H(u) + \ps{-\dive \nabla u}{h}_\X + O(\nr{h}^2).
\end{align*}
D'où
\[
  \boxed{\nabla J_H(u) = -\dive(\nabla u) = -\Delta u.}
\]

\subsection*{Fonction $J_{PM}(u) = \sum_{x} \sqrt{\nr{\nabla u(x)}^2 + 1}$}

\paragraph{Convexité.}
Posons $\psi(p) = \sqrt{|p|^2 + 1}$ pour $p\in\R^2$. On a $\psi(p) = \nr{(p,1)}_{\R^3}$ : c'est la norme
euclidienne composée avec l'application affine $p\mapsto(p,1)$, donc $\psi$ est convexe.
Pour $t\in[0,1]$ :
\[
  \psi(tp+(1-t)q) = \nr{t(p,1) + (1-t)(q,1)} \le t\,\psi(p) + (1-t)\,\psi(q).
\]
Ainsi $J_{PM} = \sum_x \psi\circ(\nabla u)(x)$ est une somme de fonctions convexes composées avec une application
linéaire : $J_{PM}$ est \textbf{convexe}. Comme pour $J_H$, $J_{PM}(u+c) = J_{PM}(u)$ pour toute constante $c$,
donc $J_{PM}$ \textbf{n'est pas strictement convexe}.

\paragraph{Gradient.}
$\psi$ est différentiable sur $\R^2$, de gradient $\nabla\psi(p) = \dfrac{p}{\sqrt{|p|^2+1}}$. Donc
\begin{align*}
  J_{PM}(u+h) &= J_{PM}(u) + \sum_x \frac{\nabla u(x)}{\sqrt{\nr{\nabla u(x)}^2+1}}\cdot \nabla h(x) + o(\nr{h}) \\
  &= J_{PM}(u) + \ps{-\dive\!\left(\frac{\nabla u}{\sqrt{\nr{\nabla u}^2+1}}\right)}{h}_\X + o(\nr{h}),
\end{align*}
\[
  \boxed{\nabla J_{PM}(u) = -\dive\!\left(\frac{\nabla u}{\sqrt{\nr{\nabla u}^2+1}}\right).}
\]
On retrouve l'EDP de Perona--Malik (1.3) avec $g(\xi) = 1/\sqrt{\xi^2+1}$.

\begin{remarque}
Pour la TV lissée $J^{\epsilon}_{TV}(u) = \sum_x \sqrt{\nr{\nabla u(x)}^2+\epsilon}$, la hessienne de
$\psi_\epsilon(p)=\sqrt{|p|^2+\epsilon}$ vaut
$\dfrac{(|p|^2+\epsilon)\,I - p\,p^{T}}{(|p|^2+\epsilon)^{3/2}}$, de plus grande valeur propre
$1/\sqrt{|p|^2+\epsilon} \le 1/\sqrt{\epsilon}$. Comme $\nr{\nabla}^2 \le 8$, on obtient
$\nabla J^{\epsilon}_{TV}$ lipschitzien de constante $L \le 8/\sqrt{\epsilon}$, qui explose quand $\epsilon\to 0$ :
c'est pourquoi le pas $\tau$ doit être pris petit (de l'ordre de $\sqrt{\epsilon}/8$) pour la TV.
\end{remarque}

\section*{Question 2 -- Stricte convexité et gradient de $J_D$}

On considère $J_D(u) = \frac12\sum_x |u(x)-f(x)|^2 = \frac12\nr{u-f}_\X^2$.

\paragraph{Stricte convexité.}
Avec la même identité que pour $J_H$ (appliquée à $a = u-f$, $b = v-f$) :
\[
  J_D(tu+(1-t)v) = t J_D(u) + (1-t) J_D(v) - \tfrac12\,t(1-t)\nr{u-v}_\X^2 .
\]
Pour $u\neq v$ et $t\in\,]0,1[$, le dernier terme est strictement positif, donc l'inégalité est stricte :
$J_D$ est \textbf{strictement convexe}. Son unique minimiseur est $u=f$.

\paragraph{Gradient.}
\[
  J_D(u+h) = \tfrac12\nr{u-f+h}_\X^2 = J_D(u) + \ps{u-f}{h}_\X + \tfrac12\nr{h}_\X^2,
  \qquad\text{donc}\qquad
  \boxed{\nabla J_D(u) = u - f.}
\]

\paragraph{Conséquence.}
$J_\lambda = \lambda J_D + J_R$ est strictement convexe dès que $\lambda>0$ (somme d'une fonction strictement
convexe et d'une fonction convexe) : le minimiseur $u_\lambda$ est unique.

\section*{Question 7 (bonus) -- Relations (1.13) et (1.14)}

On utilise le laplacien centré avec conditions périodiques :
\[
  (\Delta f)(k,l) = f(k+1,l) + f(k-1,l) + f(k,l+1) + f(k,l-1) - 4 f(k,l),
\]
les indices étant pris modulo $m$ et $n$.

\paragraph{Propriété de translation.}
Si $g(k,l) = f(k+a,l+b)$ (indices modulo $m,n$), le changement d'indice $k' = k+a$, $l'=l+b$ dans (1.11) donne
\[
  \F(g)(p,q) = e^{\,j\frac{2\pi}{m}pa}\, e^{\,j\frac{2\pi}{n}qb}\, \F(f)(p,q).
\]

\paragraph{Relation (1.13).}
Par linéarité de la TFD :
\begin{align*}
  \F(\Delta f)(p,q) &= \F(f)(p,q)\left( e^{j\frac{2\pi p}{m}} + e^{-j\frac{2\pi p}{m}} + e^{j\frac{2\pi q}{n}} + e^{-j\frac{2\pi q}{n}} - 4 \right) \\
  &= \F(f)(p,q)\left( 2\cos\tfrac{2\pi p}{m} + 2\cos\tfrac{2\pi q}{n} - 4 \right).
\end{align*}
Avec $1-\cos(2a) = 2\sin^2(a)$ :
\[
  2\cos\tfrac{2\pi p}{m} - 2 = -4\sin^2\tfrac{\pi p}{m}, \qquad 2\cos\tfrac{2\pi q}{n} - 2 = -4\sin^2\tfrac{\pi q}{n},
\]
d'où
\[
  \boxed{\F(\Delta f)(p,q) = -4\,\F(f)(p,q)\left(\sin^2\tfrac{\pi p}{m} + \sin^2\tfrac{\pi q}{n}\right).} \tag{1.13}
\]

\paragraph{Relation (1.14).}
On applique la TFD à la condition d'optimalité (1.10) $\lambda(u-f) - \Delta u = 0$ :
\[
  \lambda\,\F(u) + 4\left(\sin^2\tfrac{\pi p}{m} + \sin^2\tfrac{\pi q}{n}\right)\F(u) = \lambda\,\F(f),
\]
\[
  \boxed{\F(u)(p,q) = \frac{\lambda\,\F(f)(p,q)}{\lambda + 4\left(\sin^2\frac{\pi p}{m} + \sin^2\frac{\pi q}{n}\right)}.} \tag{1.14}
\]
Pour $\lambda>0$ le dénominateur est strictement positif, donc la solution est bien définie et unique.
Interprétation : c'est un \textbf{filtre passe-bas}. En $(p,q)=(0,0)$ le gain vaut $1$ (la moyenne de $f$ est
conservée) ; les hautes fréquences sont atténuées d'autant plus que $\lambda$ est petit.

\section*{Question 8 -- Gradient de $J_A$, calcul de $G^T$ et $M^T$}

\paragraph{Gradient de $J_A(u) = \frac12\nr{Au-f}_\X^2$.}
\begin{align*}
  J_A(u+h) &= \tfrac12\nr{Au - f + Ah}^2 = J_A(u) + \ps{Au-f}{Ah}_\X + \tfrac12\nr{Ah}^2 \\
  &= J_A(u) + \ps{A^{T}(Au-f)}{h}_\X + O(\nr{h}^2),
\end{align*}
\[
  \boxed{\nabla J_A(u) = A^{T}(Au - f).}
\]
(L'énoncé écrit $A^T(A(u-f))$ ; la bonne expression est $A^T(Au-f)$. Les deux coïncident pour le masque, voir plus bas.)

\paragraph{Convolution : calcul de $G^T$.}
La convolution s'écrit $(G*u)(x_k) = \sum_l G(x_k - x_l)\,u(x_l)$, donc la matrice associée a pour coefficients
$G_{kl} = G(x_k - x_l)$. Sa transposée vérifie
\[
  (G^{T})_{kl} = G_{lk} = G(x_l - x_k) = \tilde G(x_k - x_l), \qquad \tilde G(x) := G(-x).
\]
$G^T$ est donc la convolution avec le noyau retourné $\tilde G$. Si $G$ est symétrique ($G(-x)=G(x)$, ce qui est le cas
d'un noyau isotrope comme le noyau gaussien), alors $\boxed{G^T = G}$ et
\[
  \nabla J_G(u) = G*(G*u - f). \tag{2.16}
\]
(Aux bords, cette égalité n'est exacte que pour des conditions périodiques ; avec un prolongement symétrique,
elle est approchée.)

\paragraph{Inpainting : calcul de $M^T$.}
$M$ est la matrice diagonale avec $M_{kk} = 1$ si $x_k\in\mathcal M$ et $0$ sinon. Elle est diagonale donc
$\boxed{M^T = M}$, et c'est un projecteur : $M^2 = M$. Comme $f$ n'est observée que sur $\mathcal M$, on a :
\[
  \nabla J_M(u) = M^T(Mu - f) = M^2 u - Mf = M(u - f) = (u-f)\,M,
\]
c'est-à-dire $\nabla J_M(u)(x) = u(x)-f(x)$ si $x\in\mathcal M$ et $0$ sinon, ce qui est la relation (2.17).

\section*{Question 13 -- Propriétés des modèles et des paramètres}

\subsection*{Les modèles}

\begin{itemize}[leftmargin=*]
  \item \textbf{Déconvolution + TV lissée} (2.18).
    Le noyau gaussien est un passe-bas : $G$ a des valeurs propres très petites pour les hautes fréquences,
    donc inverser seulement le terme d'attache ($\lambda\to\infty$) amplifie le bruit.
    La TV impose une image constante par morceaux : elle redonne des contours nets, mais
    efface les textures fines (effet \og cartoon \fg{} / marches d'escalier).
    La fonctionnelle est strictement convexe : $\psi_\epsilon$ est strictement convexe, et sur les constantes
    $G * c = c \neq 0$ (noyau normalisé), donc le minimiseur est unique.
    La régularisation de Tikhonov serait peu adaptée ici, car elle lisserait encore l'image déjà floue.

  \item \textbf{Inpainting + TV lissée} (2.19).
    Hors du masque, seul le terme de régularisation agit : on prolonge l'information connue vers les trous.
    La TV favorise des lignes de niveau courtes : elle prolonge les contours et conserve des transitions nettes,
    mais peut créer des contours droits artificiels dans les grandes zones manquantes.

  \item \textbf{Inpainting + Tikhonov} (2.20).
    Dans les trous, la condition d'optimalité devient $\Delta u = 0$ : on fait une interpolation
    \textbf{harmonique}, c'est-à-dire une diffusion isotrope. Elle donne de bons résultats pour les zones presque
    uniformes (ciel, eau, sable) et les trous fins (grille du perroquet), mais elle floute les contours qui
    traversent un grand trou.
\end{itemize}

Pour les deux inpaintings, la fonctionnelle est strictement convexe dès que le masque est non vide
($M\mathbf 1 \neq 0$) : la solution est unique.

\subsection*{Les paramètres}

\begin{itemize}[leftmargin=*]
  \item $\boldsymbol{\lambda}$ : compromis attache aux données / régularité.
    En déconvolution, $\lambda$ grand donne une image plus nette mais bruitée ; $\lambda$ petit donne une image
    lisse et \og cartoon \fg{}.
    En inpainting, les pixels connus ne sont pas bruités, donc on peut prendre $\lambda$ grand ; $\lambda$ n'agit
    que sur $\mathcal M$ et a peu d'influence sur le remplissage des trous.
  \item $\boldsymbol{\epsilon}$ : plus il est petit, plus $J^\epsilon_{TV}$ est proche de la vraie TV (contours plus
    nets), mais plus la constante de Lipschitz $8/\sqrt\epsilon$ est grande, ce qui force un pas $\tau$ plus petit
    et ralentit la convergence.
  \item $\boldsymbol{\tau}$ : la descente de gradient converge pour $\tau < 2/L$, avec
    $L \le \lambda\nr{A}^2 + 8$ (Tikhonov) ou $L \le \lambda\nr{A}^2 + 8/\sqrt\epsilon$ (TV), et $\nr{G}\le 1$,
    $\nr{M}=1$.
  \item $\boldsymbol{K}$ : en inpainting, l'information se propage par diffusion, donc lentement : il faut un
    grand nombre d'itérations pour remplir les grands trous (d'où $K$ de l'ordre de plusieurs milliers).
\end{itemize}

\section*{Question 14 -- Conditions d'optimalité pour $\phi_1$ et $\phi_2$}

On minimise $J_\lambda(u) = \frac{\lambda}{2}\nr{u-f}_\X^2 + \sum_x \phi\big(\nr{\nabla u(x)}\big)$.
Pour $\phi$ dérivable avec $\phi'(0)=0$, le même calcul qu'à la question 1 donne
\[
  \nabla J_\lambda(u) = \lambda(u-f) - \dive\!\left(\frac{\phi'(\nr{\nabla u})}{\nr{\nabla u}}\,\nabla u\right),
\]
d'où la condition d'optimalité (équation d'Euler--Lagrange) :
\[
  \boxed{\lambda(u-f) - \dive\!\left(\frac{\phi'(\nr{\nabla u})}{\nr{\nabla u}}\,\nabla u\right) = 0.}
\]

\begin{center}
\renewcommand{\arraystretch}{1.8}
\begin{tabular}{@{}llll@{}}
\toprule
 & $\phi(\xi)$ & $\phi'(\xi)$ & $\phi'(\xi)/\xi$ \\
\midrule
TV lissée & $\sqrt{\xi^2+1}$ & $\dfrac{\xi}{\sqrt{\xi^2+1}}$ & $\dfrac{1}{\sqrt{\xi^2+1}}$ \\
$\phi_1$ & $\dfrac{\xi^2}{1+\xi^2}$ & $\dfrac{2\xi}{(1+\xi^2)^2}$ & $\dfrac{2}{(1+\xi^2)^2}$ \\
$\phi_2$ & $\log(1+\xi^2)$ & $\dfrac{2\xi}{1+\xi^2}$ & $\dfrac{2}{1+\xi^2}$ \\
\bottomrule
\end{tabular}
\end{center}

Les conditions d'optimalité sont donc :
\begin{align*}
  \phi_1 :\quad & \lambda(u-f) - \dive\!\left(\frac{2\,\nabla u}{\big(1+\nr{\nabla u}^2\big)^2}\right) = 0, \\
  \phi_2 :\quad & \lambda(u-f) - \dive\!\left(\frac{2\,\nabla u}{1+\nr{\nabla u}^2}\right) = 0 .
\end{align*}

\paragraph{Détail des dérivées.}
Pour $\phi_1(\xi) = \dfrac{\xi^2}{1+\xi^2} = 1 - \dfrac{1}{1+\xi^2}$, on a $\phi_1'(\xi) = \dfrac{2\xi}{(1+\xi^2)^2}$.
Pour $\phi_2(\xi)=\log(1+\xi^2)$, on a $\phi_2'(\xi) = \dfrac{2\xi}{1+\xi^2}$.

\paragraph{Non-convexité.}
$\phi_1''(\xi) = \dfrac{2(1-3\xi^2)}{(1+\xi^2)^3}$ est négatif pour $\xi > 1/\sqrt3$, et
$\phi_2''(\xi) = \dfrac{2(1-\xi^2)}{(1+\xi^2)^2}$ est négatif pour $\xi>1$ :
$\phi_1$ et $\phi_2$ ne sont \textbf{pas convexes}. La condition d'optimalité est alors seulement nécessaire :
la descente de gradient converge vers un point critique qui dépend de l'initialisation $u^0=f$.

\paragraph{Mise à l'échelle.}
Les niveaux de gris étant dans $[0,255]$, on utilise en pratique $\varepsilon^2\phi(\xi/\varepsilon)$, dont
le coefficient de diffusion est $\phi'(s)/s$ avec $s = \nr{\nabla u}/\varepsilon$. Le paramètre $\varepsilon$
joue le rôle d'un seuil de contraste : un gradient bien plus grand que $\varepsilon$ est considéré comme un contour
et peu diffusé. Comme $\phi'(s)/s \le 2$, la descente de gradient est stable pour $\tau < 2/(\lambda+16)$.

\end{document}
